\documentclass[11pt,a4paper]{article}
\usepackage[margin=24mm]{geometry}
\usepackage{amsmath,amssymb,amsthm,hyperref}
\newtheorem{lemma}{Lemma}
\newtheorem{proposition}[lemma]{Proposition}
\newtheorem{theorem}[lemma]{Theorem}
\newcommand{\gen}[1]{\langle #1\rangle}
\newcommand{\Cent}{\mathrm C}
\newcommand{\Norm}{\mathrm N}
\newcommand{\Aut}{\operatorname{Aut}}
\title{Small centralizers and the rank-four obstruction for Ree groups}
\author{Human Compression project: recovery of L039}
\date{}
\begin{document}\maketitle
This paper supplies a replacement for the rank-four argument of
Leemans, Schulte and Van Maldeghem~\cite{LSV}. It uses their classical
subgroup setting and Pierro's centralizer-intersection theorem, rather
than the normality assertion in their final rank-four argument.
The previously reviewed companion \emph{Involution blocks and generation
in the small Ree groups} supplies the complete rank-three construction,
the exclusion of proper almost simple extensions, and the upper bound
of four. Its full proof remains an acquired part of the recovery.
Neither the subgroup classifications nor the theory of Ree groups are
being reproved here.

\section{Elementary obstruction lemmas}
All actions of automorphisms are on the right: $x^{\alpha\beta}
=(x^\alpha)^\beta$, and $x^{\operatorname{inn}(u)}=u^{-1}xu$.
A string marking $(r_0,\ldots,r_{d-1})$ consists of independent
involutions with $[r_i,r_j]=1$ for $|i-j|>1$ and the intersection
property
\[
 \gen{r_i:i\in I}\cap\gen{r_i:i\in J}
 =\gen{r_i:i\in I\cap J}.
\]
Independence follows from this property and the requirement that every
generator is nonidentity. In a nonabelian simple group, adjacent
generators cannot commute: a commuting adjacent pair would cut the
string into two nontrivial commuting normal subgroups.

\begin{lemma}[Dihedral product obstruction]\label{dihedral}
Let $E$ be an elementary abelian $2$-group and let $D_{2n}$ be dihedral
of order $2n$, with $n$ odd. No subgroup of $E\times D_{2n}$ has a
rank-three string marking $(x,y,z)$ with both $xy$ and $yz$ of order
greater than two.
\end{lemma}
\begin{proof}
An involution projects either to the identity or to a reflection in
$D_{2n}$. A trivial projection would make it central, so all three
projections are reflections. Since $x,z$ commute, their reflections
coincide: two distinct reflections in an odd dihedral group do not
commute. Write $xy=(e,r)$ and $yz=(f,r^{-1})$, where $r$ has odd order
$d>1$. Squaring removes the central $2$-coordinates, and $r^2$ generates
$\gen r$. Thus both $\gen{x,y}$ and $\gen{y,z}$ contain
$\{1\}\times\gen r$. This nontrivial odd-order subgroup cannot be
contained in $\gen y$, contradicting the intersection property.
\end{proof}

\begin{lemma}[An end centralizer]\label{end}
Let a centerless group $G$ be generated by a string $(a,b,c,d)$ of
involutions, with $c,d$ noncommuting and $\Cent_G(c,d)$ abelian.
If $K=\gen{b,c,d}$ has an involution in its centralizer in $G$, then
such a marking is impossible.
\end{lemma}
\begin{proof}
If $t$ is that involution, then $a,t\in\Cent_G(c,d)$, so $a$ commutes
with $t$. The other three generators already commute with $t$.
Consequently $t\in Z(G)$, contrary to the hypothesis.
\end{proof}

\begin{lemma}[The field-degree reduction]\label{degree}
Let $G=\gen{a,b,c,d}$ be as above, and suppose
$\Cent_G(c,d)$ is a subgroup of a Klein four-group. If a nonidentity
automorphism $\tau$ of odd order fixes $b,c,d$, then $|\tau|=3$.
Moreover $\Cent_G(c,d)$ is a Klein four-group on whose nonidentity
elements $\tau$ acts cyclically.
\end{lemma}
\begin{proof}
The centralizer is $\tau$-invariant and contains $a$. If $\tau$ fixed
$a$, it would fix every generator, hence be the identity. Its induced
action is therefore a nontrivial odd-order automorphism of a subgroup
of a Klein four-group. The centralizer must be the full four-group,
and the induced action has order three. Hence $\tau^3$ fixes $a$ and
$b,c,d$, so is the identity on $G$. This proves both assertions.
\end{proof}

\begin{theorem}[Two conjugate fixed subgroups]\label{fixed}
Let $G$ be a finite centerless group with elementary abelian Sylow
$2$-subgroups. Assume that the centralizer of each pair of noncommuting
involutions is either elementary abelian of characteristic three or
a subgroup of a Klein four-group. There is no rank-four string
marking $(a,b,c,d)$ with noncommuting adjacent pairs for which
\[
 K=\gen{b,c,d}=\operatorname{Fix}(\tau),\qquad
 L=\gen{a,b,c}=\operatorname{Fix}(\sigma),
\]
where $\tau$ has order three and $\sigma$ has the same image as $\tau$
in $\Aut(G)/\operatorname{Inn}(G)$.
\end{theorem}
\begin{proof}
There is a unique $u\in G$ with $\sigma=\tau\operatorname{inn}(u)$.
Both automorphisms fix $b,c$, hence $u\in A:=\Cent_G(b,c)$.
If $u=1$ the two fixed subgroups coincide, contrary to independence;
thus $u\ne1$. The group $A$ is $\tau$-invariant because $\tau$ fixes
$b,c$. The intersection property gives $K\cap L=D:=\gen{b,c}$.

First suppose $A$ is elementary abelian of characteristic three.
The nontrivial group $U=\gen{u,u^\tau,u^{\tau^2}}$ is $\tau$-invariant.
An automorphism of order dividing three of a nontrivial finite
$3$-group has a nonidentity fixed point: its orbits have lengths one
or three, so the number of fixed points is divisible by three.
Choose $1\ne v\in U$ fixed by $\tau$. Since $U\le A$ is abelian,
$v$ commutes with $u$, and is therefore fixed by $\sigma$ as well.
Thus $v\in K\cap L=D$, while $v\in A=\Cent_G(D)$. This puts an
element of order three in $Z(D)$, impossible for a dihedral group,
whose center has order at most two.

Otherwise $A$ is a $2$-group of exponent two, so $u^2=1$.
Since $a^\sigma=a$, we have $a^\tau=a^u$. The elements $a,a^\tau$
are distinct because $a\notin K$. Both centralize $c,d$, and hence
commute. Their four-group $V=\gen{a,a^\tau}$ is normalized by $u$,
which interchanges its two displayed generators. Consequently
$\gen{V,u}$ is a nonabelian $2$-group of order eight. This contradicts
the elementary abelian Sylow $2$-subgroups of $G$.
\end{proof}

\section{The retained Ree structure}
Write $R(q)={}^2G_2(q)$ for $q=3^f$, $f$ odd, and first assume $q\ge27$.
The following inputs are retained explicitly from the classical
description summarized in~\cite[Section 2.2]{LSV} and
\cite[Section 2 and Lemma 3.13]{Pierro}.
\begin{enumerate}
\item The group $G=R(q)$ is centerless and simple, its Sylow
$2$-subgroups are elementary abelian of order eight, and
$\Aut(G)/G$ is cyclic of order $f$.
\item Its maximal subgroups are proper subfield Ree groups,
point stabilizers of order $q^3(q-1)$, involution centralizers
$C_2\times\operatorname{PSL}_2(q)$, four-group normalizers
$(C_2^2\times D_{(q+1)/2}):C_3$, and the two cyclic Hall normalizers
$C_{q+1\pm\sqrt{3q}}:C_6$. In the four-group normalizer the indicated
$C_3$ acts nontrivially on the four-group. Here the subscript of $D$
always denotes its order.
\item For noncommuting involutions $x,y$, the group $\Cent_G(x,y)$
is elementary abelian of characteristic three or a subgroup of a
Klein four-group. This is precisely the noncommuting case of
Pierro's Lemma 3.13; the commuting case has a larger centralizer.
\item A standard subfield subgroup $R(s)$, where $s=3^h$ and $h\mid f$,
is the fixed subgroup of the corresponding field automorphism of
order $f/h$. Every subgroup of this subfield type is conjugate to
the standard one. This fixed-field identification is part of the
standard matrix realization of the Ree groups; see Kleidman
\cite[Lemma 6.2]{Kleidman}, as explicitly restated in
\cite[Section 6, p.~7]{Nagy}. The automorphism acts by the indicated
field map on matrix entries. For $s=3$ its fixed
subgroup is $R(3)=\operatorname{PSL}_2(8):C_3$, whose derived group is
$\operatorname{PSL}_2(8)$.
\item The group $\operatorname{PSL}_2(8)$ has maximal subgroups
$C_2^3:C_7$, $D_{14}$ and $D_{18}$. Its nonidentity element orders
are $2,3,7,9$. Its split tori of order seven and their dihedral
normalizers form single conjugacy classes. In $R(3)$ the outer
element acts as the field automorphism of $\operatorname{PSL}_2(8)$.
These are the usual degree-eight instance of the subgroup and
automorphism descriptions for $\operatorname{PSL}_2$ used in~\cite{LSV}.
\item Every element of order three or nine fixes exactly one point
in the natural action of $R(q)$ of degree $q^3+1$. The point
stabilizer has $2$-part two. These facts are recorded in
\cite[Tables 1 and 3]{Pierro}.
\end{enumerate}
These inputs do not include the disputed normality assertion or the
middle-dihedral assertion of source Lemma 3.13. No finite computation
is a premise of the arguments below.

\begin{proposition}[Possible end parabolics]\label{parabolic}
In a hypothetical rank-four string marking of $G=R(q)$, each end
parabolic is a subfield $R(s)$ with $s\ge27$, or the derived group
$\operatorname{PSL}_2(8)$ of a subfield $R(3)$.
\end{proposition}
\begin{proof}
Consider $K=\gen{b,c,d}$. It is proper, has the intersection property,
and its adjacent pairs do not commute. Choose a maximal subgroup
containing it. Point stabilizers and cyclic Hall normalizers have
Sylow $2$-subgroups of order two, whereas $\gen{b,d}$ has order four.
An involution centralizer is impossible by Lemma~\ref{end}.
In a four-group normalizer, every involution lies in the normal
subgroup $C_2^2\times D_{(q+1)/2}$ of index three. Write
$(q+1)/2=2n$; since $q\equiv3\pmod8$, $n$ is odd.
Lemma~\ref{dihedral} excludes this case.

Only a subfield Ree group remains. Repeat this argument within a
proper subfield as long as $K$ is proper there. An involution
centralizing $K$ at any stage still lies in $G$, so the same exclusion
applies. The strictly decreasing subfield orders make the process
terminate, either with $K=R(s)$, $s\ge27$, or with $K\le R(3)$.
In the latter case all its involutions lie in
$R(3)'=\operatorname{PSL}_2(8)$. A proper subgroup of this last group
lies in one of its listed maximal subgroups. The involutions in
$C_2^3:C_7$ all lie in the elementary abelian normal subgroup and
cannot give a noncommuting adjacent pair. The odd dihedral groups
$D_{14}$ and $D_{18}$ cannot contain two distinct commuting
involutions. Hence $K=\operatorname{PSL}_2(8)$. The dual argument
applies to $L=\gen{a,b,c}$.
\end{proof}

\begin{proposition}[Reduction to a cubic field extension]\label{cubic}
If an end parabolic is $R(s)$, $s\ge27$, then $q=s^3$.
If it is $\operatorname{PSL}_2(8)$, then $q=27$.
\end{proposition}
\begin{proof}
Choose the corresponding field automorphism $\tau$ fixing that end
parabolic pointwise. In the first case its fixed subgroup is $R(s)$;
in the second it is a conjugate of $R(3)$ and its order is $f$.
This automorphism is nonidentity because its fixed subgroup is proper
in the simple group $G$. The centralizer of the last adjacent pair
contains the first generator, an involution, so the characteristic-three
alternative is impossible and it is a subgroup of a four-group.
Lemma~\ref{degree} gives $|\tau|=3$, proving the two conclusions.
\end{proof}

\section{The two final cases}
If $q>27$, Proposition~\ref{cubic} excludes an end parabolic of type
$\operatorname{PSL}_2(8)$. Both ends are therefore $R(s)$ with
$q=s^3$, and hence are conjugate. Let $\tau$ be the field automorphism
fixing one end. Conjugating $\tau$ by an inner automorphism gives an
automorphism $\sigma$ whose fixed subgroup is the other end and whose
image in $\Aut(G)/G$ is the same. Theorem~\ref{fixed} gives a
contradiction. This argument replaces all three parts of source
Lemma 3.13 and its use in Lemma 3.14; in particular it does not assume
the asserted restriction on the middle dihedral order.

It remains to treat $q=27$. Both end parabolics must now be
$\operatorname{PSL}_2(8)$. We give the ordered-triple argument in
the full automorphism group, where the required transitivity holds.

\begin{lemma}\label{seven}
In this exceptional case, $|cd|=|ab|=7$.
\end{lemma}
\begin{proof}
The possible orders are $3,7,9$. If $cd$ had order three or nine,
its unique fixed point in the natural Ree action would be fixed by
both $c$ and $a$: conjugation by $c$ inverts $cd$, whereas $a$
centralizes it. The four-group $\gen{a,c}$ would then lie in a point
stabilizer of $2$-part two, a contradiction. Thus $|cd|=7$.
The dual argument gives $|ab|=7$.
\end{proof}

\begin{proposition}[The exceptional automorphism obstruction]\label{exception}
The group $R(27)$ has no rank-four string marking.
\end{proposition}
\begin{proof}
For an ordered triple $(x,y,z)$ occurring here, $D=\gen{y,z}$ is
$D_{14}$ inside a subfield derived group $K\cong\operatorname{PSL}_2(8)$,
and $1\ne x\in\Cent_G(D)$. All such subfield derived groups are
conjugate; their $D_{14}$ subgroups are conjugate inside them.
Moreover the normalizer of $D$ in $R(3)$ induces the full automorphism
group of $D_{14}$. To see the last assertion explicitly, use the
split torus $T=\gen r$ of order seven in $\operatorname{SL}_2(8)$
and its basis-swapping involution $w$. Inner conjugation by $T$ and
$w$ induces the translations and inversion on the seven reflections
$wr^i$. The field map on $\mathbb F_8$ induces multiplication by two
on their indices. Since $\gen{-1,2}=\mathbb F_7^\times$, these actions
give all affine permutations $i\mapsto h+ki$ and are transitive on
ordered pairs of distinct reflections. This element of order three
in $R(3)$ is separate from the following outer automorphism of $G$.

Fixing such an ordered pair $(y,z)$, let $\tau$ be the order-three
field automorphism of $G$ fixing its containing $R(3)$ pointwise.
No nonidentity involution of $\Cent_G(D)$ is fixed by $\tau$:
any such involution would belong to $R(3)$, hence to $K$, whereas
$\Cent_K(D)=1$. The latter equality is immediate in the displayed
$2$-dimensional model: an element centralizing the split torus is
diagonal and commutation with $w$ forces it to be the identity.
The centralizer is a subgroup of a four-group and contains $x$;
it is therefore the full four-group, and $\tau$ cycles its three
nonidentity elements. Combining these actions proves transitivity
of $\Aut(G)$ on the relevant ordered triples.

Apply this to $(a,c,d)$ and $(d,b,a)$, which satisfy the same
conditions by Lemma~\ref{seven}. Some automorphism interchanges
$a$ and $d$, and thus induces a transposition on the nonidentity
elements of $V=\gen{a,d}$. However $\Norm_G(V)$ induces only $C_3$
on $V$, by the retained four-group-normalizer structure. In
$\Aut(G)$ the quotient $\Norm_{\Aut(G)}(V)/\Norm_G(V)$ has odd
order, since it embeds in $\Aut(G)/G$. Its induced group on $V$
therefore has odd order as well and cannot contain a transposition.
This is a contradiction.
\end{proof}

\begin{theorem}[Complete principal conclusion]\label{principal}
For a simple small Ree group $R(q)$ with $q\ge27$, the only rank
of a string C-group representation is three. Every such $R(q)$ has
a representation of rank three. No proper almost simple extension
of $R(q)$ is generated by involutions, and $R(3)$ is not generated
by involutions.
\end{theorem}
\begin{proof}
The propositions and Theorem~\ref{fixed} exclude rank four.
The previously accepted companion proves the bound of four,
the existence in rank three, and both odd-quotient exclusions.
Ranks one and two generate cyclic or dihedral groups and are
impossible for a nonabelian simple group. Together these statements
give exactly the principal conclusion of~\cite[Theorem 1.1]{LSV}.
\end{proof}

The source does not enumerate all rank-three markings, and neither
does this recovery. The theorem is a recovery of the published
existence and exclusion statements; no new historical priority,
conjecture answer, complete formalization or numerical factor is claimed.

\section{Auxiliary claims and the subgroup diagram}
The preceding proof deliberately avoids most normalizer calculations
in the source. We record their precise status so that this change of
route does not hide a source claim. In particular, an implication
whose hypothesis is a rank-four Ree marking is discharged by the
nonexistence theorem; it is not thereby asserted as an unconditional
subgroup-normalizer theorem.

\paragraph{Dihedral and projective-line calculations.}
In $D_{2n}=\gen{r,w}$, a subgroup of order $2m$, $m>1$ and $m\mid n$,
is $\gen{r^{n/m},wr^j}$. A rotation $r^i$ normalizes it exactly when
$n/m\mid2i$; multiplying by its reflection gives the other half of
the normalizer. The normalizer therefore has order $2m$ if $n/m$
is odd and $4m$ if it is even. This proves preprint Lemma 1,
published Lemma 2.1.

For $L=\operatorname{PSL}_2(q)$, $q=3^f$ with $f$ odd, we retain
the standard split and nonsplit torus description: their orders are
$(q-1)/2$ and $(q+1)/2$, their normalizers are dihedral of twice
those orders, and a nonidentity semisimple element of order greater
than two has its torus as centralizer. Nontrivial unipotents have
order three. A dihedral subgroup $D_{2d}$, $d\ge3$, cannot have
$d=3$: conjugation inverting a nontrivial unipotent, after its unique
projective fixed point is moved to infinity, would require an upper
triangular determinant-one matrix whose diagonal ratio is $-1$.
That ratio is a square, whereas $-1$ is not a square for
$q\equiv3\pmod4$. Thus its rotation lies in a torus, and
$2d\mid q-1$ or $q+1$. In particular $4\nmid d$, and even $d$ can
only occur in the latter case. Its normalizer lies in the normalizer
of its rotation torus. The preceding dihedral calculation determines
it. This is the complete calculation needed in source Lemma 3.5.

For clarity about Lemma 3.6, let $J=D_{2d}\le C_2\times L$ with
$d\ge3$. Project onto $L$, obtaining a nonabelian dihedral group $D$.
Either $J=C_2\times D$, or $J$ is the graph of a homomorphism
$\chi:D\to C_2$. In the graph case,
\[
 N_{C_2\times L}(J)
 =C_2\times\{g\in N_L(D):\chi(x^g)=\chi(x)\text{ for all }x\in D\}.
\]
If the rotation order of $D$ is odd, its only possible nontrivial
homomorphism to $C_2$ distinguishes rotations and reflections and is
preserved by every automorphism. If it is even, $|N_L(D):D|$ cannot
be two: that would give a subgroup of $L$ with order divisible by
eight, whereas $|L|$ has $2$-part four. Inner automorphisms preserve
$\chi$. Consequently, in every case
\[
 N_{C_2\times L}(J)=C_2\times N_L(D).
\]
When $|D|=2d$, the latter normalizer is $C_2\times D_{2d}$, except
when $2d\mid q+1$ and $(q+1)/(2d)$ is even, in which case it is
$C_2\times D_{4d}$. When $J=C_2\times D$, we have $|D|=d$, $d/2$
odd, and $N_L(D)=D_{2d}$, so the result is again $C_2\times D_{2d}$.
These formulas prove the embedding-independence assertion in
Lemma 3.6(d). The ``according as'' order in its part (a) interchanges
the two signs; the containing torus normalizer has order $q-1$
when $2d\mid q-1$, and $q+1$ when $2d\mid q+1$.

Source Lemma 3.4 is immediate from the intersection property:
$r_0$ centralizes $G_{01}$, but if it normalized the proper nontrivial
$G_0$ then $G_0$ would be normal in the simple generated group.
The preceding arguments prove the end-parabolic form actually needed
from Lemmas 3.7 and 3.10 directly, with its noncommuting-adjacency
hypothesis explicit. To recover the nondegenerate string version of
the independent subgroup statement in Lemma 3.7, one may stop the
descent at $C_2\times\operatorname{PSL}_2(s)$ instead of excluding
it. Dickson's retained subgroup list at odd powers of three consists
of subfield groups, odd-order point stabilizers, the two torus
normalizers, and the $\operatorname{PSL}_2(3)=A_4$ endpoint.
The first three reductions work as above. In $C_2\times A_4$, all
involutions generate an elementary abelian group, so cannot supply
the required adjacent pairs. Finally a subgroup projecting onto a
simple $\operatorname{PSL}_2(s)$ is either that factor or its direct
product with $C_2$, since a simple nonabelian group has no nontrivial
homomorphism to $C_2$. This gives the four source families under the
nondegenerate string convention used in its proof. The literal claim
without that convention is too broad: $C_2^3\le R(q)$ itself has a
rank-three string C-group marking, with adjacent labels two.

For source Lemma 3.8, all involutions of $R(3)$ lie in
$K=\operatorname{PSL}_2(8)$. The product of two distinct involutions
has order $3,7$ or $9$. Each corresponding dihedral group $D$ has
odd rotation order and is self-normalizing in $K$: use the split
torus of order seven or the nonsplit torus of order nine and the
dihedral normalizer calculation (including its order-six subgroup).
Thus every involution normalizing $D$ already lies in $D$, which
is exactly the asserted $N^0_{R(3)}(D)=D$.

\paragraph{The conjugacy-class diagram.}
Figure 1 is interpreted as in the source: its boxes are specific
conjugacy classes, not every subgroup of the abstract isomorphism
type. It applies when $3\mid f$, so that a subfield $H=R(27)$ exists.
Put $J=R(3)$, $K=J'$, and let $D=D_{14}$ be a dihedral subgroup
of $K$. Put $N=N_G(D)$. We retain the standard subfield conjugacy
and normalizer data used for this diagram:
\[
 N_G(H)=H,\quad N_G(J)=J,\quad N_G(K)=J,\quad
 N=(C_2^2\times D):C_3,\quad |N|=168.
\]
The subfield and normal-dihedral data are recorded in
\cite[Lemmas 2.5, 2.7, 2.9--2.10]{Pierro}; the $K=J'$ normalizer
is also part of the source's subfield description.
The group $N$ contains four subgroups abstractly isomorphic to
$D_{14}$, but only one belongs to the displayed class, namely its
normal factor $D$. Indeed the other three are obtained by multiplying
its reflections by a fixed nonidentity element of the four-group.
An element of order three cycles these three twists; their
normalizers have order $56$, rather than $168$, so they are not
$G$-conjugate to $D$. To compute those normalizers, note that any
normalizing element preserves the characteristic rotation subgroup
of order seven. In its torus normalizer a rotation must square into
that subgroup, and the $C_3$-factor cannot preserve a nonzero
four-group twist. This gives exactly $C_2^2\times D_{14}$.
The unique normal $D_{14}$ is characteristic in $N$, whence
$N_G(N)=N_G(D)=N$.

All subfield and normal-dihedral classes in question have the same
fusion in $H$ as specified above. Also $|J|=1512$, $|K|=504$, and
$N_K(D)=D$, so $K$ contains $504/14=36$ members of the displayed
dihedral class. The class sizes and all edge multiplicities in the
figure now follow by double counting incident subgroup pairs:
\[
\begin{array}{c|c|c}
 \text{upper and lower classes}&\text{lower per upper}&\text{upper per lower}\\\hline
 G,H&|G|/|H|&1\\
 H,J&|H|/1512&1\\
 J,K&1&1\\
 K,D&36&4\\
 H,N&|H|/168&1\\
 N,D&1&1
\end{array}
\]
The five proper class sizes are respectively $|G|/|H|$,
$|G|/1512$, $|G|/1512$, $|G|/168$, and $|G|/168$.
For example the number of overgroups $K$ of a fixed $D$ is
$(|G|/1512)36/(|G|/168)=4$. No enumeration of all subgroup
classes is being claimed. The rank-four replacement does not depend
on the uniqueness assertions read from this diagram.

\paragraph{Edition and correction boundary.}
The published eighteen-page primary was inspected through its
publisher-hosted PDF text representation; its exact PDF bytes were
not available to the local downloader. The locally bound preprint
remains the byte-pinned primary. Published Lemmas 3.12--3.14 retain
the proof steps at issue. In Lemma 3.14, the asserted normality of
$H=D:C_3$ inside $N=(C_2^2\times D):C_3$ would imply that
$H/D=C_3$ is normal in $N/D=A_4$, which it is not.
The old normality inference is therefore not a premise here.
Likewise, for a fixed ordered generating pair of $D_{14}$, the
inner stabilizer is its four-group centralizer and cannot cycle
that centralizer's three nonidentity elements. Proposition~\ref{exception}
uses the full automorphism group instead. Finally, the last
generation sentence of source Lemma 3.15 must use the third
involution outside the torus normalizer, as proved in the companion,
not the first involution already chosen inside it.

The source's question about chiral polytopes of higher rank is kept
as an unanswered source-recorded question. We make no current-status
or priority claim for that question. All source implications assuming
a rank-four Ree marking are settled by Theorem~\ref{principal};
their intermediate normalizer reasoning is superseded, rather than
silently certified.

\begin{thebibliography}{9}
\bibitem{LSV} D. Leemans, E. Schulte and H. Van Maldeghem,
\emph{Groups of Ree type in characteristic 3 acting on polytopes},
Ars Math. Contemp. \textbf{14} (2018), 209--226,
\url{https://cage.ugent.be/~hvm/artikels/235.pdf}.
The local preprint is arXiv:1501.01125v1, eighteen pages;
its Lemmas 12--14 correspond to published Lemmas 3.12--3.14.
\bibitem{Pierro} E. Pierro, \emph{The M\"obius function of the small Ree groups},
Australas. J. Combin. \textbf{66}(2) (2016), 142--176,
\url{https://ajc.maths.uq.edu.au/pdf/66/ajc_v66_p142.pdf}.
\bibitem{Kleidman} P. B. Kleidman, \emph{The maximal subgroups of the
Chevalley groups $G_2(q)$ with $q$ odd, the Ree groups ${}^2G_2(q)$,
and their automorphism groups}, J. Algebra \textbf{117} (1988), 30--71.
\bibitem{Nagy} G. P. Nagy, \emph{Embeddings of Ree unitals in a projective
plane over a field}, Finite Fields Appl. \textbf{74} (2021), 101875.
The field-centralizer statement is explicitly restated in Section~6,
p.~7 of arXiv:2007.10464v3,
\url{https://real.mtak.hu/156937/1/2007.10464.pdf}.
\end{thebibliography}
\end{document}
